\documentclass[11pt,a4paper]{article}
\usepackage[T1]{fontenc}
\usepackage{lmodern}
\usepackage[margin=27mm]{geometry}
\usepackage{amsmath,amssymb,amsthm,booktabs,microtype}
\usepackage[hidelinks]{hyperref}
\newtheorem{theorem}{Theorem}
\newtheorem{proposition}[theorem]{Proposition}
\newtheorem{lemma}[theorem]{Lemma}
\theoremstyle{definition}
\newtheorem*{desiredtheorem}{Desired theorem (unproved)}
\theoremstyle{remark}
\newtheorem{remark}[theorem]{Remark}
\newcommand{\Ostar}{O^{*}}
\newcommand{\Pcal}{\mathcal P}
\newcommand{\E}{\mathbb E}
\newcommand{\Prb}{\mathbb P}
\emergencystretch=2em
\clubpenalty=10000
\widowpenalty=10000
\title{A Learning and Coverage Dichotomy for Max 2-SAT\\
\large Two-sided comparisons, conditional dominance, and shared constraints}
\author{}
\date{7 October 2026}
\begin{document}
\maketitle

\begin{abstract}
We investigate how overlapping learned constraints might yield a
density-independent exponential saving for Max 2-SAT. Fractional packing,
strengthened by conditioning on shared variables, gives a computable
completion bound compatible with branching. We prove the required
inheritance inequalities and construct logarithmic-width clauses with
linear joint information. Two-sided comparisons lead to an entropy
dichotomy: nearly implied comparisons give either immediate pruning or
information already present. Joint infeasibility certificates exploit one
case, including balanced variables with identical comparison profiles.
Small exact experiments show both useful pruning and substantial overhead.
A family with almost $2^n$ strict local optima identifies a remaining
obstruction. A precise desired frontier theorem combines certified pruning
with structural solutions; it remains unproved. No fixed-base algorithm
for growing density is established.
\end{abstract}

\section{What the dichotomy should accomplish}

The target is a deterministic polynomial-space algorithm taking
$\Ostar((2-\delta)^n)$ time for a fixed $\delta>0$, independently of a
growing density bound. We use the canonical unweighted $B=1$ setting of the
earlier notes unless stated otherwise. A bound on average clause density is
not a bound on maximum variable degree.

The useful alternatives are not merely ``large neighborhood'' and
``dominated variable''. They are:
\begin{enumerate}
\item Large residual neighborhoods permit shared exposure: a small selected
set of variables affects many targets, potentially creating useful gaps.
\item Smaller residual neighborhoods permit learning; accumulated learned
constraints can then be exploited through their joint restrictions.
\end{enumerate}
The second alternative permits learned widths tending to infinity. The
Kulikov--Kutzkov learning method motivates this alternative~\cite{KK}.
The present note studies the accounting and exploitation of already sound
constraints. It does not establish a new sound rule for learning arbitrary
clauses, nor a theorem that the useful constraint families must arise.

At a node, all hard constraints must be valid for that node's search
responsibility. Constraints justified only in one subtree must not be
copied unconditionally into its sibling. The certificates below can ignore
some valid constraints, but never introduce an unjustified one.

\paragraph{Exposure is available.}
For reference, suppose a neighborhood hypergraph has $n$ vertices,
minimum edge size $s$, maximum incidence degree $d$, and $1\le k\le s$.
There is a $k$-transversal of size at most
\[
 \min\left\{n,\frac n s\bigl(3k+8\ln(d+1)\bigr)\right\}.
\]
To prove this, select each vertex with probability
$p=(2k+8\ln(d+1))/s$ if $p<1$, and repair deficient edges by adding at most
$k$ vertices each. A Chernoff bound makes each edge deficient with
probability at most $1/(d+1)$, and there are at most $nd/s$ edges.
The displayed bound follows by expectation. If $p\ge1$, select all vertices.
Thus in a $d$-regular graph with $d\to\infty$, $k=o(d)$ exposures per
neighborhood cost $o(n)$ selected variables. This is an unsigned statement.
Under uniform assignments, $t$ exposed neighbors change one target's
effective gap by a signed sum with standard deviation $\sqrt t$, not by
$t$ deterministically. Covering does not settle the signed analysis.

\section{Exact information and two barriers}

Let $R$ be the set of $r$ free variables and $\Pcal$ the active hard
constraints. Put
\[
 Z(\Pcal)=\#\{x\in\{0,1\}^R:x\models\Pcal\},\qquad
 I(\Pcal)=r-\log_2 Z(\Pcal).
\]
An inconsistent family has $I=+\infty$. Adding a constraint $C$ to a
nonempty family changes the information by
\[
 I(\Pcal\cup\{C\})-I(\Pcal)
 =-\log_2\Prb[C\mid\Pcal],
\]
where the conditional distribution is uniform on the satisfying assignments.
The quantity is joint information, not the sum of unconditional clause
contributions. Computing it exactly is generally difficult.

\begin{proposition}[Residual-width barrier]\label{prop:width}
For $m$ ordinary non-tautological clauses, each on distinct variables and
with residual length at least $w$, if $m2^{-w}<1$, then
\[
 Z\ge2^r(1-m2^{-w}),\qquad
 I\le-\log_2(1-m2^{-w}).
\]
\end{proposition}
\begin{proof}
A length-$\ell$ clause is false on a $2^{-\ell}$ fraction of the Boolean
cube. Apply the union bound to all clause-violation events. Overlap and
literal signs do not invalidate that bound.
\end{proof}

Consequently, polynomially many ordinary clauses all of length
$\omega(\log n)$ carry $o(1)$ bits of joint information at that state.
This does not exclude savings later, once clauses become shorter, or from
succinct cardinality constraints representing many ordinary clauses.
It also does not exclude logarithmic or slower-growing learned widths.

\begin{proposition}[Exposure alone preserves assignment mass]
For a fixed ordinary clause $C$ of length $\ell$, assign any $t$ of its
variables uniformly. It remains unsatisfied with probability $2^{-t}$;
conditional on this, its residual length is $\ell-t$. Hence its expected
remaining violation probability is $2^{-\ell}$. More generally, for any
fixed hard family and variable $x$,
\[
 Z(\Pcal)=Z(\Pcal\mid x=0)+Z(\Pcal\mid x=1).
\]
\end{proposition}

Shortening exposes existing restrictions and can make them easier to
detect. It does not by itself create new exact information after the cost
of all branches is included. A running-time proof must credit new learning,
forced assignments, stronger certificates, or cheaper solution of the
remaining objective. In particular, analysis must account for both the
branches satisfying a learned clause and those shortening it.

\section{A fractional-packing completion certificate}

For each active constraint $C_i$, let $E_i\subseteq R$ be a support on
which it depends and let
\[
 p_i=2^{-|E_i|}\#\{x\in\{0,1\}^{E_i}:x\models C_i\},\qquad
 c_i=-\log_2p_i.
\]
Ignore identically true constraints. An individually false constraint
immediately refutes the node. For the remaining constraints, choose
nonnegative weights satisfying
\begin{equation}
 \sum_{i:v\in E_i}w_i\le1\quad(v\in R).             \label{eq:packing}
\end{equation}
Define
\begin{equation}
 W_w=\sum_iw_ic_i,\qquad
 U_w=2^r\prod_i p_i^{w_i}=2^{r-W_w}.                  \label{eq:bound}
\end{equation}
Zero weights omit their factors. A weight-one factor is allowed for an
identically false constraint and makes $U=0$.

\begin{theorem}[Completion bound and inheritance]\label{thm:packing}
For every feasible weight vector, $Z(\Pcal)\le U_w$. If a variable is
branched on, retaining the same weights in the two restricted children
gives
\begin{equation}
 U_{w,0}+U_{w,1}\le U_w.                             \label{eq:split}
\end{equation}
Weights may be reoptimized in each child, provided the inherited witness
remains an option. Newly learned constraints receive inherited weight zero.
The same inequality then continues to hold.
\end{theorem}
\begin{proof}
Let $J$ be the constraints containing the branching variable $x$.
For $i\in J$, $p_i=(p_{i,0}+p_{i,1})/2$. Other probabilities are unchanged.
By~\eqref{eq:packing}, $\sum_{i\in J}w_i\le1$, so the weighted geometric
mean $g(z)=\prod_{i\in J}z_i^{w_i}$ is concave on the nonnegative orthant.
Consequently
\[
 \frac12\prod_{i\in J}p_{i,0}^{w_i}
 +\frac12\prod_{i\in J}p_{i,1}^{w_i}
 \le\prod_{i\in J}\left(\frac{p_{i,0}+p_{i,1}}2\right)^{w_i}.
\]
Multiplying by $2^r$ and the unchanged factors proves~\eqref{eq:split}.
Feasibility persists after removing an assigned variable. Repeat the
inequality to complete assignments. Every satisfying assignment has leaf
bound one; other assignments have a nonnegative bound. Summing proves
$Z\le U$. Further valid restrictions and minimization over candidates
including the inherited one can only decrease child bounds.
\end{proof}

The static bound is a product-space inequality of Finner/Shearer type;
it is not a new probability inequality. Read-$k$ versions and their
limitations are discussed in~\cite{GLSS}. The point here is the explicit
branching-compatible witness and its extension to conditioning below.

\paragraph{Computability.}
For an ordinary clause of residual width $\ell$,
$p_i=1-2^{-\ell}$. For an unweighted cardinality constraint permitting
at most $b$ true literals among $\ell$ distinct variables,
\[
 p_i=2^{-\ell}\sum_{j=0}^{\min(b,\ell)}\binom\ell j,
\]
with $p_i=0$ if $b<0$. Numerically polynomially bounded integer-weight
constraints can use a knapsack dynamic program. General constraints do
not automatically have efficiently computable probabilities.

For a fixed family, maximizing $W_w$ subject to~\eqref{eq:packing} is a
linear program with coefficients $c_i$. Exact optimality is unnecessary:
any rational feasible vector supplies a valid witness. Numerical LP
outputs must be repaired or certified for feasibility before use.

\paragraph{How much ordinary packing can certify.}
If every variable appears in at most $D$ active supports, $w_i=1/D$ is
feasible and gives $I\ge D^{-1}\sum_i c_i$. Here $D$ is the incidence
degree of the \emph{learned} family; it need not equal the original graph's
degree. This is a useful sufficient bound but can deteriorate with $D$.

There is a stronger limitation on the certificate itself. If
$c_i\le\eta|E_i|$ for every factor, then
\begin{equation}
 W_w\le\eta\sum_iw_i|E_i|\le\eta r.                 \label{eq:efficiency}
\end{equation}
For ordinary clauses all of width at least $w$, one can take
\[
 \eta=\frac{-\log_2(1-2^{-w})}{w}=O\!\left(\frac1{w2^w}\right).
\]
Thus fractional packing of long ordinary clauses alone cannot certify a
linear information gain, regardless of how many clauses there are.
For a cardinality constraint with allowance $\beta\ell$, $\beta<1/2$
fixed, its information can instead be $\Theta(\ell)$. The framework
therefore treats ordinary clauses and strong cardinality clauses differently.

\section{Conditioning on shared variables}

Choose a set $H\subseteq R$, with $h=|H|$. It will be called a core;
this means only a chosen set of shared variables, not a graph separator.
For every assignment $a\in\{0,1\}^H$, compute the residual probabilities
$p_i(a)$ and choose weights $w_i(a)$ subject to
\[
 \sum_{i:v\in E_i\setminus H}w_i(a)\le1
 \qquad(v\in R\setminus H).
\]
The choices may differ between core assignments. Write
$W(a)=\sum_iw_i(a)(-\log_2p_i(a))$. A locally contradicted assignment
has $W(a)=+\infty$. Define
\begin{equation}
 U_H=2^{r-h}\sum_{a\in\{0,1\}^H}2^{-W(a)},\qquad
 I_H=r-\log_2U_H
     =-\log_2\E_a[2^{-W(a)}].                         \label{eq:core}
\end{equation}

\begin{theorem}[Conditional certificate]\label{thm:core}
We have $Z(\Pcal)\le U_H$. Store $H$ and its assignment-indexed weight
vectors as the witness. Restrict that witness in each child and retain it
among the candidates. Then
\[
 U_{H,0}+U_{H,1}\le U_H.
\]
For $h=O(\log n)$, polynomially many constraints with polynomial-time
local counting permit polynomial-size witnesses and polynomial work per
certificate evaluation, up to the numerical issue discussed below.
\end{theorem}
\begin{proof}
Apply Theorem~\ref{thm:packing} independently under each fixed core
assignment and sum. Branching on $x\in H$ partitions the sum exactly.
For $x\notin H$, apply~\eqref{eq:split} separately to each summand.
The number of summands is $2^h$. New cores and stronger weight vectors may
be tried, while retaining the restricted old witness guarantees inheritance.
\end{proof}

Enlarging the core never worsens the best certificate: reuse each old
weight vector in both extensions and apply the same concavity argument.
Feasible weights can improve further because variables moved into the core
no longer impose packing capacities.

\paragraph{Numerical implementation.}
Unlike the earlier integer weighted-sum certificate, $U_H$ need not be
an integer. For rigorous search, compute an upper approximation
$\widetilde U\in[U,(1+\epsilon)U]$, with exact handling of zero factors.
Use rational feasible weights and directed error bounds for logarithms,
exponentials, and summation. Polynomial precision suffices when the family,
weight representations, and local probability representations have
polynomial size. Inheritance gives
$\widetilde U_0+\widetilde U_1\le(1+\epsilon)\widetilde U$.
For $\epsilon=n^{-2}$, the inflation over at most $n$ branching levels is
bounded by a constant. Cut only when the certified upper bound is less
than one. The prototype below checks mathematical inequalities numerically;
it does not implement directed rounding and is not a certified solver.

\paragraph{Rare core assignments cannot be ignored.}
For every $a$ with finite $W(a)$,
\begin{equation}
 I_H\le h+W(a).                                      \label{eq:worstcore}
\end{equation}
This follows by retaining only its summand in~\eqref{eq:core}.
If $h=O(\log n)$, one core assignment with $W(a)=o(n)$ prevents this
certificate from giving $\Omega(n)$ information. A high average of $W(a)$
is insufficient. Such surviving branches need further learning or a
different structural argument. This is an explicit place where the
proposed dichotomy must continue recursively.

\section{Long clauses can jointly be very strong}

Let $Y=\{y_1,\ldots,y_q\}$ and $X=\{x_1,\ldots,x_t\}$.
For each $a\in\{0,1\}^q$, let $D_a(Y)$ be the ordinary clause falsified
exactly by the pattern $a$.

\begin{proposition}[A linear-size family with logarithmic width]
The family
\begin{equation}
 \Pcal=
 \{D_a(Y):a\ne0^q\}
 \ \cup\ 
 \{x_i\lor y_1\lor\cdots\lor y_q:1\le i\le t\}       \label{eq:gate}
\end{equation}
has exactly one satisfying assignment. It has $2^q-1+t$ clauses of
width $q$ or $q+1$. Taking $t=2^q$ gives $n=t+q$, $O(n)$ clauses,
and every width $\Theta(\log n)$. Nevertheless its joint information is
exactly $n$ bits.
\end{proposition}
\begin{proof}
The first group forces $Y=0^q$. The second group then forces every
$x_i=1$. Conditioning on $H=Y$ detects a false core clause on every
nonzero pattern; at the zero pattern, disjoint unit factors on $X$ have
weights one. Hence the conditional certificate also gives $U_H=1$.
\end{proof}

Without conditioning, all supports contain $Y$, so the packing weights
sum to at most one. The optimal unconditioned information is only
$-\log_2(1-2^{-q})$. This separates the conditional construction from
ordinary fractional packing by a linear amount.

A second family, containing $x_i\lor D_a(Y)$ for every $i,a$, is equivalent
to $\bigwedge_i x_i$. Its core is free, and conditioning on $Y$ gives
$U_H=2^q$, again exactly. Both examples are abstract hard-constraint
families. No claim is made that the Max 2-SAT learning rules generate
them with the required frequency or cost.

\paragraph{Growing degree does not prevent the example.}
Take disjoint copies of~\eqref{eq:gate}, each with $t=2^q$ targets.
The learned hypergraph has maximum incidence degree $d=2t-1$, widths
$q$ or $q+1=\Theta(\log d)$, and $O(n)$ clauses. It still has a unique
satisfying assignment. With $q\to\infty$ this gives examples at growing
degree, not only when the degree is linear in the total number of variables.
The union of all local cores has size
$nq/(t+q)=O(n\log d/d)=o(n)$, while each individual core is logarithmic.
The example also has many low-degree target variables; no regularity claim
is intended.

\begin{proposition}[Composing disjoint local certificates]
Suppose valid subfamilies have pairwise disjoint variable supports $B_j$.
For each, let $U_j$ be a completion upper bound satisfying the inherited
branching inequality, possibly obtained using its own core. Then
\[
 U=2^{r-|\bigcup_jB_j|}\prod_j U_j
\]
is a completion upper bound for their conjunction and satisfies the same
branching inequality. Other active constraints may be ignored or used in
additional candidate certificates.
\end{proposition}
\begin{proof}
The subfamily completion counts multiply on disjoint supports. A branch
in $B_j$ changes only its factor; an uncovered-variable branch splits its
free factor exactly in half. Apply the local branching inequality.
\end{proof}

Thus many local cores can be used without enumerating their entire union
at once. In the disjoint-copy example every local bound equals one, so
their product certifies the unique global assignment with polynomial work.
This suggests looking for a packing of useful local cores as well as one
globally shared core. It does not require the full Max 2-SAT objective to
decompose: the bound concerns the selected hard constraints only.

At the other extreme, if every ordinary learned clause contains the same
positive literal $z$, then at least $2^{r-1}$ assignments satisfy the whole
family. Thus $I\le1$, however many distinct clauses accumulate.
Branching on $z$ exposes a potentially constrained child and a child in
which the learned family disappears. The latter child's original Max
2-SAT objective may still be arbitrary. Clause redundancy alone therefore
does not establish useful structure in that objective.

\section{What covering and duality provide}

For ordinary learned clauses, a useful core-selection objective is to meet
selected supports $E_i$ in at least $|E_i|-b$ variables. Under any core
assignment, each selected clause is either satisfied or becomes a clause
of width at most $b$. This is a multicover problem with nonuniform demands.
Signs decide which clauses survive, so a support-only objective must be
followed by the conditional sum in~\eqref{eq:core}.

The fractional-packing LP also supplies an exact, but quantitatively
limited, covering alternative. Its dual is
\begin{equation}
 \min\sum_{v\in R}y_v
 \quad\text{subject to}\quad
 y_v\ge0,\quad \sum_{v\in E_i}y_v\ge c_i.             \label{eq:dual}
\end{equation}
We consider nonconstant factors with $0<p_i<1$, so their supports are
nonempty and the usual $w_i\le1$ restrictions are redundant.

\begin{proposition}[Packing or a hitting set]\label{prop:dual}
Suppose a selected group of factors has $|E_i|\le b$ and $c_i\ge a>0$.
Let $W^*$ be the optimal packing value for the full family. There is a set
meeting every support in the selected group, of size at most $bW^*/a$.
\end{proposition}
\begin{proof}
Take an optimal dual vector and put $H=\{v:y_v\ge a/b\}$.
If a selected support missed $H$, its at most $b$ coordinates would
sum to less than $a$, contradicting its dual constraint. Also
$|H|a/b\le\sum_vy_v=W^*$.
\end{proof}

This gives a genuine packing-or-covering statement: either $W^*$ is large,
or all selected supports meet a set bounded in terms of $W^*$.
But it hits each support only once, and for long ordinary clauses
$a$ can be about $2^{-b}$. Its bound then becomes too weak. It does not
establish a logarithmic core, near-complete coverage, or a small graph
separator. A stronger theorem must use clause provenance, repeated
shortening, or interactions between opposite literals.

\paragraph{Stronger factors.}
The factors need not be individual clauses. Any consequence with an
efficiently countable satisfying set may be used. Two useful choices are:
\begin{itemize}
\item a conjunction of constraints whose union of supports has
$O(\log n)$ variables, counted by enumeration;
\item a nonnegative weighted sum of cardinality inequalities, counted by
the integer dynamic program from the earlier branching note.
\end{itemize}
These can be inserted as additional factors or maintained as separate
candidate certificates. Inherited witnesses are retained in either case.
The approaches are complementary: two opposite cardinality inequalities
may sum to a contradiction that marginal packing misses completely.
Conversely, resolving sign patterns through a shared core can reveal
restrictions hidden from a single weighted-sum relaxation.

\section{Two-sided comparisons and conditional dominance}
\label{sec:comparison}

The question ``does a fresh constraint add information?'' is incomplete
if it refers only to the second child. A comparison that is redundant
there may exclude the corresponding part of the first child.

\begin{proposition}[Partition by the better value]\label{prop:partition}
Write an assignment as $(x,y)$, let $f_b(y)=f(b,y)$, and suppose the hard
constraint $C(y)$ does not involve $x$. Define the exact comparison
\[
 c_x(y)\quad\Longleftrightarrow\quad f_1(y)>f_0(y).
\]
With ties assigned to $x=0$,
\[
 \max_{x,y:C(y)} f(x,y)
 =\max\left\{
     \max_{y:C(y)\wedge\neg c_x(y)} f_0(y),\quad
     \max_{y:C(y)\wedge c_x(y)} f_1(y)
   \right\}.                                      \label{eq:partition}
\]
An empty maximization has value $-\infty$.
\end{proposition}
\begin{proof}
For each fixed feasible $y$, the displayed partition keeps exactly the
better of its two extensions, with the stated tie convention.
\end{proof}

Consequently, $C\Rightarrow c_x$ eliminates the entire $x=0$ child.
If the implication holds only under a partial assignment $\alpha$,
then $C\wedge\alpha\Rightarrow c_x$ eliminates that portion of the child.
When the implication is only approximate, the first child needs to search
only the exceptional set $C\wedge\neg c_x$. This gives a concrete meaning
to exploiting constraints that are already present.

For Max 2-SAT, $f_1-f_0$ is an affine function of the other Boolean
variables, so the exact comparison is a threshold constraint. In the
earlier notes' notation, the preferred literal has flip difference
$L-p$: its child keeps $L\ge p$, and its opposite child keeps
$L\le p-1$. The asymmetric weak/strict thresholds assign every tie
to the preferred literal. A learned clause may be a weaker consequence
of this threshold. Merely knowing that such a necessary consequence
holds does \emph{not} prove which value is better. A sufficient improvement
certificate can eliminate its certified region; an exact comparison gives
the full partition~\eqref{eq:partition}.

\paragraph{Feasibility and scope are part of the rule.}
If genuine hard constraints involve $x$, put $C_b(y)=C(b,y)$.
The correct two retained domains are
\begin{align*}
 &C_0(y)\wedge\bigl(\neg C_1(y)\vee f_0(y)\ge f_1(y)\bigr),\\
 &C_1(y)\wedge\bigl(\neg C_0(y)\vee f_1(y)>f_0(y)\bigr).
\end{align*}
For example, $C_0\Rightarrow C_1\wedge c_x$ suffices to discard the first
child. The implication $C\wedge(x=1)\Rightarrow c_x$ alone does not:
it can express a relation between $x$ and its neighbors rather than
a globally preferred value of $x$.

Optimization-preserving pruning constraints need a further distinction:
they describe the node's search responsibility, not necessarily the full
original feasible domain. Transferring a solution to the other value of
$x$ must respect the established search invariant, possibly because a
dominating solution is covered elsewhere. The prototype below avoids a
new transfer assumption: it forms the children of the existing sound
two-sided search and only refutes conjunctions within those children.
Likewise, a certificate available only after solving the first child
does not retroactively remove the cost of solving it. Prospective
comparison certificates are needed to save that work.

\subsection{Disagreement mass measures both children}

Let $\Omega\subseteq\{0,1\}^R$ be the nonempty current hard domain,
with $r=|R|$, and let $A$ be uniform on $\Omega$.
Let $b_x(y)$ be the comparison-preferred value of $x$.
For a genuine hard domain it is feasibility-aware: choose $1$ precisely
when $C_1(y)$ holds and either $C_0(y)$ fails or $f_1(y)>f_0(y)$.
Otherwise choose $0$ on the projection of $\Omega$.
For a pruning domain, use a comparison whose filtering is justified by
the search invariant. Define
\[
 \beta_x=\Prb\bigl[A_x\ne b_x(A_{-x})\bigr].
\]
Filtering both children by this comparison leaves exactly
$(1-\beta_x)|\Omega|$ assignments in total. Its exact information gain is
\begin{equation}
 g_x=-\log_2(1-\beta_x).                         \label{eq:disagreement}
\end{equation}
Thus a comparison can be useful even when its second-child constraint
adds no new information there. If $C$ is independent of $x$, every
feasible $y$ has two extensions and $\beta_x=1/2$. This does not yield
one bit of gain at every later step: accumulated constraints can couple
$x$ to $y$, making $\beta_x$ very small.

\section{A proved near-implication entropy dichotomy}
\label{sec:entropy}

Let $h_2(q)=-q\log_2q-(1-q)\log_2(1-q)$ denote binary entropy, with
the usual endpoint convention. The following is a counting statement;
it does not assume that the conditional probabilities are efficiently
computable.

\begin{theorem}[Nearly fixed comparisons: pruning or existing information]
\label{thm:entropy}
Let $\Omega\subseteq\{0,1\}^R$ be nonempty, $A$ uniform on $\Omega$,
and $T\subseteq R$ nonempty. For every $x\in T$, suppose there is a
constant $a_x\in\{0,1\}$ such that
\[
 \Prb\bigl[b_x(A_{-x})\ne a_x\bigr]\le\eta.
\]
For $\varepsilon>0$, $\eta\ge0$, and
$\eta+\varepsilon<1/2$, at least one of the following holds:
\begin{enumerate}
\item Some $x\in T$ has $\beta_x\ge\varepsilon$, so its two-sided
comparison removes at least an $\varepsilon$ fraction of $\Omega$.
\item The existing domain already satisfies
\begin{equation}
 \log_2|\Omega|
 \le r-|T|\bigl(1-h_2(\eta+\varepsilon)\bigr).
                                                        \label{eq:entropy}
\end{equation}
\end{enumerate}
\end{theorem}
\begin{proof}
If the first alternative fails, then for every $x\in T$,
\[
 \{A_x\ne a_x\}\subseteq
 \{A_x\ne b_x(A_{-x})\}\cup\{b_x(A_{-x})\ne a_x\}
\]
gives $\Prb[A_x\ne a_x]\le\varepsilon+\eta$.
Subadditivity of entropy and monotonicity of $h_2$ on $[0,1/2]$ give
\[
 \log_2|\Omega|=H(A)
 \le\sum_{x\in R}H(A_x)
 \le r-|T|+|T|h_2(\eta+\varepsilon).
\]
\end{proof}

For example, with $\eta=\varepsilon=1/8$, either one comparison removes
at least $12.5\%$ of the current domain, or that domain already contains
at least $0.18872|T|$ bits of information. If $|T|\ge\tau r$ for a
fixed $\tau>0$, the second alternative is linear in $r$.
The theorem formalizes the proposed dichotomy when many comparison
outcomes are almost constant under the old constraints. No degree gap
assumption is needed for the statement.

There are two remaining gaps between this result and an algorithmic
running-time bound. First, true conditional probabilities and model counts
are not supplied by the theorem. An upper bound on the exception count
divided by an upper bound on the total count is not, in general, an upper
bound on the exception probability. Efficient certificates must justify
the required implication, mass reduction, or completion bound directly.
Second, small $\beta_x$ need not mean that $b_x$ is almost constant.
The old constraints may already encode balanced, nonconstant response
relations. That case is not covered by the theorem's hypothesis and can
have very little information, as Section~\ref{sec:stable} shows.

\section{Exploiting implication before opening a child}
\label{sec:joint}

\subsection{Identical comparison profiles without initial imbalance}

Suppose $T$ is an independent set of variables in the interaction graph,
$Y=R\setminus T$, and all centers $x\in T$ have the same difference
\[
 f(x=1,y)-f(x=0,y)=h(y),\qquad y\in\{0,1\}^Y.
\]
The difference is independent of the other centers because there are no
clauses between them. Use a common tie convention. In the unconstrained
soft optimization problem, or whenever the relevant transfers are sound,
the partition $h(y)>0$ versus $h(y)\le0$ simultaneously determines a
best value for every center. The optimum is attained with all centers
equal to that common best response.

After the first center's comparison is imposed, the opposite comparison
for every later center is inconsistent. For integral $h$, the two
inequalities $h\ge1$ and $h\le0$ sum to a contradiction. Thus joint
constraints can force the remaining centers even if each inequality
alone has no unit consequence and every initial literal gap is zero.
The construction in Section~\ref{sec:experiments} demonstrates this effect.
This is a mechanism example, not a new general algorithm for independent
sets: a large independent set is already useful structure. Its relevance
is that repeated comparison information can be exploited before a large
literal imbalance appears.

\subsection{A checkable joint-infeasibility certificate}

After assigning a candidate value and applying the existing propagation
rules, write the residual child constraints as
\[
 \sum_j c_{ij}^{z_j}\le k_i \quad (i=1,\ldots,m),
 \qquad c_{ij}^0,c_{ij}^1\ge0.
\]
Here $c_{ij}^b$ is the contribution of variable $j$ to row $i$ when
$z_j=b$; these superscripts are labels, not powers. For weights $w_i\ge0$
define
\begin{equation}
 M(w)=\sum_j\min_{b\in\{0,1\}}\sum_i w_i c_{ij}^b
                  -\sum_i w_i k_i.                    \label{eq:margin}
\end{equation}

\begin{proposition}[Joint refutation]\label{prop:refutation}
If $M(w)>0$, the child constraint system is infeasible.
\end{proposition}
\begin{proof}
Every assignment has weighted left-hand side at least
$\sum_j\min_b\sum_iw_ic_{ij}^b$, whereas satisfying all rows would make
it at most $\sum_iw_ik_i$.
\end{proof}

For example, refuting $C\wedge\neg c_x$ proves $C\Rightarrow c_x$.
In the implemented search we instead test the complete, correctly scoped
constraint system of each proposed child. One refuted child fixes the
other value; two refuted children cut the parent. The same certificate
can be used under partial assignments to cover portions of an exceptional
first branch. This is the zero-completion case of the weighted-sum
method, used prospectively for the two-sided comparison.

The prototype first tests equal weights on all rows and on every pair
of rows. If these fail, it solves the normalized linear program
\begin{align*}
 \text{maximize }&\sum_j t_j-\sum_iw_i k_i,\\
 \text{subject to }&t_j\le\sum_iw_i c_{ij}^b
                     \quad(j,\ b\in\{0,1\}),\\
 &w_i\ge0,\qquad\sum_iw_i=1.
\end{align*}
A floating-point solution only proposes multipliers. They are converted
to nonnegative rationals, and~\eqref{eq:margin} is verified exactly before
any cut. An unverified proposal is ignored. The certificate also refutes
the continuous box relaxation, so it is incomplete for Boolean
infeasibility. Its use does not claim a new linear-programming duality
principle. Polynomially many rows and polynomial certificate bit lengths
permit polynomial verification and polynomial-size linear programs.

\section{The balanced-relation obstruction}
\label{sec:stable}

The remaining case is not just ``large, balanced neighborhoods without
new clauses.'' It can be a domain in which the comparison relations are
already encoded, yet almost all of the exponential search space remains.
The following example uses $B=2$, with a complementary pair of clauses
on each edge; it is not a $B=1$ lower-bound construction.

\begin{proposition}[Almost $2^r$ strict local optima]\label{prop:stable}
For even $r$, place the two clauses
$(x_i\vee x_j)$ and $(\neg x_i\vee\neg x_j)$ on every pair $i<j$.
There are exactly $\binom r{r/2}=\Theta(2^r/\sqrt r)$ strict
one-variable local maxima, all globally optimal. On their domain every
$\beta_x$ is zero and every variable has marginal $1/2$.
\end{proposition}
\begin{proof}
If $k$ variables are one, the number of satisfied clauses is
\[
 f(x)=\binom r2+k(r-k).
\]
It is maximized at $k=r/2$. Flipping any one variable there decreases
the objective by one. At any other cardinality, a flip toward $r/2$
increases it. Hence the strict local maxima are precisely the balanced
layer. Each coordinate is balanced on that layer by complementation,
and no assignment disagrees with its unconstrained best response.
\end{proof}

Consequently, retaining all assignments satisfying every one-variable
best-response relation can leave information only
$r-\log_2\binom r{r/2}=\tfrac12\log_2 r+O(1)$.
This also gives a growing-degree example: for odd $d\to\infty$ with
$d+1$ dividing $r$, take disjoint copies of $K_{d+1}$ with the same
encoding. Interaction degree is $d$, binary clause density is $d$, and
the number of surviving strict local maxima is
\[
 \binom{d+1}{(d+1)/2}^{\,r/(d+1)}
   =2^{\,r-\Theta(r\log d/d)}.
\]
Thus pure elimination of assignments that fail a one-variable local
optimality test cannot by itself guarantee a fixed exponential saving
for growing degree. This is not a lower bound on branching algorithms:
the example is easy to optimize by the aggregate cardinality $k$, and
its components are explicit. It explains why the desired theorem must
permit structural solutions, aggregation, or additional symmetry-breaking
rules. A small separator is one possible form of structure, not a
necessary conclusion of redundant comparisons.

\section{A concrete strategy and the desired global theorem}

A possible implementation maintains the original formula, valid active
learned constraints, and the best inherited completion witness.
For candidate branching variables:
\begin{enumerate}
\item Form both children with the established learning and propagation
rules, preserving their scope and tie-breaking conventions.
\item Test their joint constraints for certified infeasibility as in
Section~\ref{sec:joint}, before exploring either subtree.
\item Evaluate the inherited witness, weighted-sum candidates, fractional
packings, and a polynomial-size list of cores of size at most
$h=c\log_2 n$ for fixed $c$.
\item Select the variable minimizing $U_0+U_1$, or maximizing
$\log_2 U-\log_2(U_0+U_1)$.
\end{enumerate}
Candidate cores can come from high incidence, the dual weights in
\eqref{eq:dual}, or a greedy multicover routine aimed at short residual
clauses. Greedily extending a core by trying each next variable is
polynomial for fixed $c$, because each trial evaluates at most $2^h$
assignments. Searching all $n^h$ cores would not be polynomial.

\paragraph{Why this is a running-time measure, not only a model count.}
For exact witnesses, every kept node has $U\ge1$, and total child mass
is at most parent mass. A subtree rooted at mass $U$ has at most
$O(nU)$ kept nodes under inherited certificates. Discarded children and
polynomial lookahead add only polynomial work. The relative-rounding
implementation changes this by a bounded multiplicative factor.

At an unconstrained root $U=2^n$, so this observation alone gives no
improved base. Let $G_t$ be accumulated certified log-decrease under
the normalized walk taking child $i$ with probability
$U_i/(U_0+U_1)$. As proved in the earlier branching note~\cite{Frontier},
a sufficient condition is, for fixed $\alpha,\gamma,\xi>0$,
\[
 \Prb\left[G_{\lfloor(1-\alpha)n\rfloor}<\gamma n\right]
 \le2^{-\xi n}.
\]
It yields $\Ostar(2^{(1-\min\{\alpha,\gamma,\xi\})n})$ time.
The conditional packing certificates can be used in that theorem just
like the weighted-sum certificates. The missing condition is global
availability of enough progress, not the validity of the local measure.

\subsection{Desired theorem: amortized dominance or structure}

The local results suggest three sources of progress: fresh comparisons
remove appreciable mass; old constraints certify that the mass is already
small; or balanced response relations permit a cheaper structural
solution. Theorem~\ref{thm:entropy} proves a restricted counting version
of the first two alternatives. Proposition~\ref{prop:stable} shows why
they cannot replace the third in a universal statement.

Here is a quantitative target that accounts for all branches. It is
intentionally called a \emph{desired theorem}, rather than a conjecture
already supported in its full generality. An algorithm realizing its
conditions is still missing.

\begin{desiredtheorem}
There are constants $\alpha,\gamma\in(0,1)$, a density bound
$\rho(n)\to\infty$, and a deterministic polynomial-space comparison-and-learning
algorithm for canonical unweighted Max 2-SAT of density at most $\rho(n)$
with the following property. Count binary branching decisions from the
root, and stop the initial search at depth
$t_0=\lfloor(1-\alpha)n\rfloor$, or earlier when a node can be discharged
by a certified structural routine. The initial search uses polynomial
work per node, and its boundary consists of:
\begin{enumerate}
\item live nodes $v\in\mathcal F$ with inherited completion bounds $U_v$
that support a continuation of cost $\Ostar(U_v)$;
\item terminal nodes $v\in\mathcal S$ with valid solution routines of
cost $\Ostar(B_v)$ and polynomial space.
\end{enumerate}
The routines and certificates preserve an optimum, have polynomial-size
representations, and satisfy the global frontier budget
\begin{equation}
 \sum_{v\in\mathcal F} U_v+\sum_{v\in\mathcal S} B_v
       \le 2^{(1-\gamma)n}.                         \label{eq:desired}
\end{equation}
Polynomial factors in this budget may be absorbed into $\Ostar$.
The constants are independent of $n$ and of the allowed growing density.
The initial target may restrict to $B=1$; the stronger target allows
$B=2$, with the ultimate density range $\rho(n)=\Theta(n)$.
\end{desiredtheorem}

A refuted terminal has $B_v=0$, and an ordinary evaluated leaf can have
$B_v=1$. A structural terminal's $B_v$ bounds its \emph{solution cost},
not its number of feasible assignments. This distinction is essential:
the balanced clique layer has almost $2^r$ assignments but admits a
polynomial-time solution. A demand that every terminal have a small
completion count would rule out exactly the useful structural exits.

\begin{proposition}[Consequence of the desired budget]
If the desired theorem holds, then the algorithm runs in
\[
 \Ostar\bigl(2^{(1-\kappa)n}\bigr)
   =\Ostar\bigl((2-\delta)^n\bigr),\qquad
 \kappa=\min\{\alpha,\gamma\},\quad
 \delta=2-2^{1-\kappa}>0,
\]
using polynomial space.
\end{proposition}
\begin{proof}
There are at most $O(2^{t_0})$ binary prefixes. Forced assignments and
polynomial local processing add only a polynomial factor. The total
continuation and terminal work is bounded by~\eqref{eq:desired}.
Depth-first traversal, polynomial-size active certificates, and the
assumed polynomial-space terminal routines give the space bound.
\end{proof}

The missing assertion is availability of this budget on states reached
by a specified sound learning algorithm. No uniform one-step dichotomy
for arbitrary hard constraints is claimed. A proof must exploit clause
provenance and the residual objective, account for exceptional core
assignments in~\eqref{eq:worstcore}, and pay for all children. The earlier
normalized-walk condition is one route for the completion-bound part;
it must be supplemented by structural exits when local comparisons
saturate. Neither the entropy theorem nor the experiments establish
these global requirements.

\section{Checks and their limits}
\label{sec:experiments}

\subsection{Completion certificates and shared cores}

The companion experiment uses exact Boolean enumeration for reference
counts, exact rational checks of packing feasibility, and floating-point
certificate evaluation with tolerance $10^{-7}$. It does not implement
the interval-arithmetic variant described above.

\begin{itemize}
\item On 250 random systems of ordinary/cardinality constraints on
3--10 variables, every certificate upper-bounded the exact count.
All 1,368 inherited branching checks, 250 added-constraint checks, and
223 core-extension checks passed. The largest positive numerical excess
in a branching inequality was $5.7\cdot10^{-13}$.
\item For the construction~\eqref{eq:gate} with $q=4,t=12$, there are
16 variables and 27 clauses of width 4 or 5. The exact count is one.
The certificate progresses as follows:
\end{itemize}
\begin{center}
\begin{tabular}{rrr}
\toprule
Core size & Completion upper bound & Certified bits\\
\midrule
0 & 61,440 & 0.093\\
1 & 57,344 & 0.193\\
2 & 49,152 & 0.415\\
3 & 32,768 & 1.000\\
4 & 1 & 16.000\\
\bottomrule
\end{tabular}
\end{center}
This illustrates a threshold effect: most of the benefit appears only
when the entire shared core is conditioned on.

\begin{itemize}
\item The all-patterns construction with the same $q,t$ has 192
width-5 clauses; conditioning on its four core variables gives the exact
count 16.
\item For $\sum_{j=1}^6x_j\le2$ and
$\sum_{j=1}^6(1-x_j)\le2$, unconditioned packing gives 22 completions,
whereas their sum refutes the system immediately.
\item We sampled 48 residual states from 12 random canonical Max 2-SAT
formulas on 14 variables, with edge probability $0.85$ and independent
balanced signs. States had 4--10 free variables and at least two active
nonvacuous pruning constraints, after domination and propagation.
Conditioning on the two most frequent variables improved packing in
24 states. It beat the earlier heuristic weighted-sum certificate in
4 states; the sum certificate was better in 34 states, with 10 ties.
All bounds remained above the exact intersection count.
\end{itemize}

These first experiments compare certificates, not whole-search runtime.
They support keeping both certificate families, without establishing
that logarithmic cores suffice in the worst case.

\subsection{Two-sided search with joint refutations}

Code and results are in
{\small\nolinkurl{explore_conditional_dominance.py}} and
{\small\nolinkurl{conditional_dominance_checks.json}}. The experiment compares the established
two-sided branching search, with domination and propagation, against
the same search with the joint-refutation test from
Section~\ref{sec:joint}. Both use the same variable order within each
comparison. Rejected children count as leaves in \emph{both} searches.
Floating-point linear programming proposes witnesses; exact rational
verification authorizes every additional cut.

\paragraph{Random instances.}
The 26 instances have 10, 12, or 14 variables (ten, ten, and six instances).
Edge probabilities lie in $[0.65,1]$, signs are independently balanced,
and complementary-pair probability is $0$ or $0.15$, with random units.
Both searches use maximum-degree branching:
\begin{center}
\begin{tabular}{lr}
\toprule
Quantity & Total\\
\midrule
Baseline leaves & 433\\
Leaves with joint refutations & 384\\
Additional joint cuts & 29\\
Branch decisions reduced to one surviving child & 23\\
Linear-program calls & 270\\
Baseline elapsed time (seconds) & 0.0129\\
Elapsed time with joint refutations (seconds) & 0.3941\\
\bottomrule
\end{tabular}
\end{center}
Both searches match exhaustive optima on all 26 instances. Leaves decrease
by $11.3\%$, but the prototype is about 30 times slower: the extra pruning
does not give a practical speedup here. Separately, 150 random weighted
constraint systems yielded 71 multiplier refutations, all checked by
exhaustive enumeration.

\paragraph{Identical balanced centers.}
For each of $t$ independent centers $x$, use $(x\vee y_j)$ on the first
$D/2$ common neighbors and $(\neg x\vee\neg y_j)$ on the remaining $D/2$.
Add a random signed clique on the neighbors and units to reduce their
unary fields. All centers have the same comparison profile and initial
deficits $(D/2,D/2)$, hence zero gap. Both searches process centers first.
\begin{center}
\begin{tabular}{rrrrrr}
\toprule
$n$ & $t$ & $D$ & Baseline leaves & Joint-test leaves & Fixings\\
\midrule
12 & 4 & 8 & 127 & 11 & 6\\
16 & 6 & 10 & 1,130 & 14 & 10\\
18 & 8 & 10 & 4,136 & 16 & 14\\
\bottomrule
\end{tabular}
\end{center}
One center's decision forces the others in each subtree. Exhaustive
optimization checks the first two instances; the baseline checks the
18-variable result. These illustrate conditional dominance, not a
worst-case bound.

\paragraph{A negative control.}
For the complete-graph Max-Cut encoding on ten variables, both searches
have 252 leaves; 248 linear-program calls find no new cuts. Exhaustive
enumeration on 12 variables gives the predicted 924 strict local optima
and only 2.1483 bits of information. As Proposition~\ref{prop:stable}
indicates, saturated comparisons require collective structural progress.

Local identities, counting bounds, and certificate rules are proved;
whole-search evidence is experimental. The global frontier budget and a
fixed-base algorithm for growing density remain desired results. No
comprehensive literature survey or general novelty claim is intended.

\begin{thebibliography}{9}
\small
\bibitem{KK}
A. S. Kulikov and K. Kutzkov,
\emph{New Bounds for MAX-SAT by Clause Learning}, CSR 2007,
LNCS 4649, pp. 194--204.
\href{https://doi.org/10.1007/978-3-540-74510-5_21}{Publisher DOI};
\href{https://pdfs.semanticscholar.org/45e9/3f8020eea4fea5a70c624e149e6a9b275a23.pdf}{authors' presentation consulted}.
\bibitem{GLSS}
D. Gavinsky, S. Lovett, M. Saks and S. Srinivasan,
\emph{A Tail Bound for Read-$k$ Families of Functions},
Random Structures \& Algorithms 47(1), 2015, pp. 99--108.
\href{https://arxiv.org/abs/1205.1478}{arXiv:1205.1478}.
\bibitem{Frontier}
\emph{Exploring the Branching Frontier for Max 2-SAT},
working note, 7 October 2026, Sections 3--4.
Local source: \texttt{max2sat\_branching\_frontier.tex}.
\end{thebibliography}
\end{document}
